% Page 034

\seite{34}

\margmark{Dat. bis Nov. 1903\\Organ. Mte.}%
        Witterung.\ob
~~Am 4.\ September wurde die hiesige Schule\ob
von einem Knäuel\unsicher{} getroffen.~~

\pend
\abschnitt{Ernte}
\pstart

Die Heuernte gab in diesem Jahre einen\ob
mittleren Ertrag. Da in der Zeit während des\ob
Heuens längere Zeit ausblieb, entwickelten\ob
sich Hafer und Gerste nur später. Desgleichen ist es\ob
dem Roggen ergangen. Roggen und Weizen sind\ob
hingegen ziemlich gut ausgefallen. Da noch die nassen\ob
Jahre dauern, so war die Ernte sehr schwierig, denn\ob
fast während des ganzen Sommers regnete es, wenn es\ob
nur die Nachmittage war. Kartoffeln ernteten eine\ob
mittlere Ernte. Obst gibt es im Ganzen nur wenig,\ob
da in letzteren Jahren noch nicht mehr Obstbäume\ob
hier angepflanzt waren.\ob
Während des Sommers hingen mehrfach schwere\ob
Gewitter über dem Himmelszelt nieder. Es\ob
schlug der Blitz in einen Ziegenstall ein.

\pend
\abschnitt{Schulfeiern}
\pstart

Der Geburtstag unseres Kaisers wurde\ob
von den Schulen Seffern, Schleid und Seffer\ohb
weich gemeinsam gefeiert. Die Festrede\ob
hielt Herr Lehrer Fritz aus Seffern, der\ob
unseren Kaiser als Friedensfürsten feierte.

Am 9.\ März wurde der Geburtstag\ob
des Prinzen Friedrich Wilhelm, des Kronprinzen,\ob
gefeiert.

Der Geburtstag unseres Kronprinzen konnte\ob
wegen Krankheit des Lehrers nicht gefeiert werden.

\pend
\abschnitt{Pfarrerwechsel}
\pstart

Anfang Mai schied unser allverehrter Herr\ob
Pfarrer Ettges aus der Pfarrei, weil er sich\ob
aus Gesundheitsrücksichten genötigt sah, eine\ob
leichtere Pfarrstelle anzunehmen. Nur ungern\ob
sah die Pfarrei ihn scheiden. Seine Liebe und Opfer-
\pend
